% =====================================================================
\chapter{Qualité, tests et limites connues}
\label{ch:qualite}
% =====================================================================

\section{La stratégie de test}

La suite de tests repose sur trois principes.

\subsection{Hors-ligne par défaut}

La suite tourne sans accès réseau : les sources sont remplacées par des
\emph{fixtures} capturées et des simulations. Les appels réels à Géorisques ou
à la BDNB constituent un étage de tests séparé.

Cette séparation n'est pas une commodité, c'est une nécessité : une suite qui
dépend de services tiers est lente, instable, et échoue pour des raisons qui
n'ont rien à voir avec le code. Une suite qui échoue sans raison est une suite
que l'on finit par ignorer.

\subsection{La dégradation testée aussi rigoureusement que le nominal}

C'est l'exigence \frref{22} de la spécification, et probablement la plus
structurante. Chaque source est éteinte indépendamment, dans chacun de ses
modes de défaillance : délai d'attente, erreur 5xx, charge malformée, réponse
vide. Dans tous les cas, la réponse attendue est un HTTP 200 partiel structuré,
avec les aléas affectés à \code{present=null} et une erreur renseignée. Jamais
un 500, jamais une exception non gagnée.

\subsection{Les tests d'abord}

Le workflow impose d'écrire le test avant l'implémentation : un commit
\code{test(TXXX)} qui échoue, puis un commit \code{feat(TXXX)} qui passe. Deux
règles complémentaires : les tests existants ne sont jamais modifiés pendant une
tâche de fonctionnalité, et un test jugé erroné est corrigé dans un change set
dédié.

\section{Inventaire de la suite}

\begin{table}[H]
\centering
\footnotesize
\caption{Les modules de test backend et leur objet.}
\label{tab:tests-backend}
\begin{tabularx}{\textwidth}{@{}l X@{}}
\toprule
\textbf{Module} & \textbf{Objet} \\
\midrule
\fichier{test\_post\_canonical\_contract.py} & Forme canonique de bout en bout \\
\fichier{test\_diagnostic\_record\_schema.py} & Champs interdits, optionnels
  toujours présents, énumération de résolution \\
\fichier{test\_georisques\_wfs.py} & Parseur GML, ordre des axes, proximité,
  boîte englobante \\
\fichier{test\_wfs\_integration.py} & Client dédié, absence de contamination \\
\fichier{test\_adresse\_georisques.py} & Jointure adresse $\rightarrow$ aléas \\
\fichier{test\_flood\_alea.py} & Classes TRI, absence $\neq$ panne \\
\fichier{test\_certification.py} & Suite de certification à réponse connue \\
\fichier{test\_geocoding\_hardening.py} & Arrondissements, seuil 0,4, 429 \\
\fichier{test\_api\_key\_auth.py} & 401/403, exemptions, clés hors journaux \\
\fichier{test\_result\_cache.py} & Clé de cache, précision, expiration TTL \\
\fichier{test\_upstream\_throttle.py} & Limitation de débit amont \\
\fichier{test\_budget\_guard.py} & Compteur mensuel, rejet du lot seul \\
\fichier{test\_chaos\_degradation.py} & Matrice de pannes par source \\
\fichier{test\_internal\_batch.py} & Isolation des éléments, équivalence unitaire \\
\fichier{test\_rate\_limit.py} & Limitation entrante \\
\fichier{test\_report\_agent.py} & Rapport narratif et repli gabarit \\
\fichier{test\_hydro.py} & Tracé hydrographique \\
\fichier{test\_panoramax.py} & Photographie et crédit \\
\fichier{test\_health.py} & Sondes de vie \\
\fichier{test\_latency\_report.py} & Percentiles de latence \\
\fichier{test\_api\_diagnostic\_offline.py} & Pipeline hors-ligne \\
\bottomrule
\end{tabularx}
\end{table}

Côté frontend, 21 modules de test couvrent le moteur de scénarios, la boucle
complète, les seuils, les vignettes, le mode VFX et l'isolation des modes.

\begin{table}[H]
\centering
\small
\caption{Volumétrie de la suite de tests au terme du stage.}
\label{tab:tests-vol}
\begin{tabular}{@{}l r l@{}}
\toprule
\textbf{Périmètre} & \textbf{Tests} & \textbf{Exécution} \\
\midrule
Backend (pytest, hors-ligne) & 214 & $<$ 60~s \\
Frontend (Vitest) & 143 & quelques secondes \\
\addlinespace
\emph{Total} & \emph{357} & \\
\bottomrule
\end{tabular}
\end{table}

\section{La suite de certification}

La certification \frref{11} mérite un développement à part. Pour quatre aléas —
inondation/PPRI, mouvement de terrain, feu de forêt, avalanche — des fixtures
ont été capturées pour une adresse \emph{connue à l'intérieur} d'un périmètre
et une adresse \emph{connue à l'extérieur}, chacune vérifiée contre les cartes
officielles du portail Géorisques.

L'attendu est strict :

\begin{itemize}[leftmargin=1.4em]
  \item à l'intérieur : \code{present=true} et
        \code{resolution="per-building"} ;
  \item à l'extérieur : \code{present=false}.
\end{itemize}

Ce n'est pas un test unitaire de plus, c'est la \textbf{porte de signature de
la version~1} : il prouve que le test géométrique dit vrai, pas seulement qu'il
s'exécute. La collecte des fixtures est la partie longue de la tâche.

Un script de capture est conservé dans le dépôt, sous le nom
\fichier{record\_certification\_fixtures.py}, avec les
coordonnées candidates et l'attendu pour chaque cas — par exemple la vallée du
Var à Nice comme cas exposé, et la colline de Cimiez comme cas non exposé.

\begin{honesty}
Cette suite est écrite, mais sa \emph{qualification humaine} — c'est-à-dire la
confrontation visuelle des fixtures aux cartes officielles — était encore en
attente à la fin du stage. Le code est là, les fixtures sont là, la validation
finale par un humain reste à faire. C'est une limite, et il vaut mieux
l'énoncer que la laisser découvrir.
\end{honesty}

\section{Les portes de non-régression}

Outre les tests, deux mécanismes automatiques protègent les invariants.

\subsection{Porte de recherche sur les champs interdits}

Une vérification échoue si les identifiants \code{niveau}, \code{score} ou
\code{recommandations} apparaissent dans un schéma émis. C'est la traduction
automatique de l'invariant \og pas de score \fg.

\subsection{Porte de recherche sur les surfaces supprimées}

Après chaque suppression, une vérification confirme qu'aucune référence ne
subsiste vers les modules retirés — \code{copernicus}, \code{mistral},
\code{lidar}, \code{rapport-narratif}. La règle \og pas de code dormant \fg{}
devient ainsi vérifiable mécaniquement.

\section{Le test d'équivalence du lot}

Un test mérite une mention : il vérifie que la sortie d'un élément de lot est
\emph{identique} à celle d'un appel unitaire pour la même adresse, à
l'horodatage près.

C'est la traduction de l'invariant \og le lot est N fois l'unitaire \fg. Sans ce
test, deux implémentations peuvent diverger lentement : un champ ajouté d'un
côté seulement, un arrondi différent, un ordre d'aléas qui change. Le test rend
cette dérive impossible.

\section{Les limites connues}

Un rapport de stage honnête doit énumérer ce qui ne fonctionne pas ou pas
encore. Le tableau~\ref{tab:limites} recense les limites assumées.

\begin{table}[H]
\centering
\small
\caption{Limites connues du système au terme du stage.}
\label{tab:limites}
\begin{tabularx}{\textwidth}{@{}l X@{}}
\toprule
\textbf{Limite} & \textbf{Description et statut} \\
\midrule
Cache non persistant en production &
  En serverless, le système de fichiers de la fonction est éphémère : le cache
  de rapports ne survit pas. Chaque rapport consomme un appel au modèle.
  Documenté. \\
\addlinespace
Store de lot en mémoire &
  Le magasin de lots est en mémoire, mono-processus, sans persistance. Perdu
  au redémarrage. Suffisant pour valider le schéma, pas dimensionné pour
  \og des centaines de milliers d'adresses \fg. \\
\addlinespace
Authentification simulée &
  Sans configuration Supabase, l'authentification reste en mode démonstration :
  la connexion fonctionne sans comptes réels. Non utilisable en production
  externe. \\
\addlinespace
RGA non vectoriel &
  Le retrait-gonflement des argiles reste communal. Un câblage sur une source
  vectorielle (par exemple les aléas argile du BRGM) demande un amendement
  constitutionnel. \\
\addlinespace
Certification non qualifiée &
  La suite de certification à réponse connue attend sa validation humaine
  contre les cartes officielles. \\
\addlinespace
Mode VFX incomplet &
  Les tâches \vfx{004}, \vfx{005} et \vfx{007} (couche Three.js, terrain,
  réalisateur caméra) ne sont pas livrées. \\
\addlinespace
Calibration de latence &
  Le budget de performance (p95 $\leq$ 10~s) est proposé, non encore calibré
  sur une campagne de mesures réelles. Un harness de mesure existe
  (\fichier{latency\_report.py}). \\
\addlinespace
Smoke test manuel non exécuté &
  La vérification manuelle de bout en bout du moteur de scénarios sur une
  adresse certifiée TRI avec pluie prévue n'a pas été exécutée : elle demande
  un jeton Mapbox et des conditions météo particulières. \\
\bottomrule
\end{tabularx}
\end{table}

\begin{honesty}
Le dernier point de ce tableau est le plus important à assumer. Une
fonctionnalité peut être couverte par des tests automatisés et n'avoir jamais
été vue fonctionner de bout en bout sur une donnée réelle. Ce n'est pas la même
chose. La vérification manuelle reste due.
\end{honesty}

\section{Performance}

Le budget de performance proposé par la constitution §6 est le suivant :

\begin{itemize}[leftmargin=1.4em]
  \item p95 $\leq$ 10~s pour \api{POST /diagnostic/adresse} avec toutes les
        sources saines ;
  \item le mode dégradé ne dépasse pas deux fois ce budget — un échec déclenche
        un repli, pas une tempête de réessais ;
  \item toute régression supérieure à 25~\% bloque la fusion.
\end{itemize}

Ces chiffres sont présentés comme \emph{proposés}, à calibrer par la mesure.
Un point d'attention est identifié : la résolution géométrique interroge
actuellement les couches WFS de façon séquentielle, chacune avec son délai
d'attente et sa réattente. Si les mesures dépassent le budget, la parallélisation
des requêtes de couches est prévue comme tâche conditionnelle (T017) — plutôt
que de relever silencieusement le budget.

\begin{honesty}
Relever un budget de performance pour qu'il corresponde à la mesure est une
manière élégante de ne jamais échouer. Le projet a choisi l'inverse : le budget
est un engagement, et s'il n'est pas tenu, c'est le code qui change.
\end{honesty}
